%!TEX root = ../template.tex
%% ====================================================================
%% Capitulo 2 -- Enquadramento
%% Esqueleto gerado 2026-08-13 a partir do indice v1
%% Seccoes 2.1-2.3 preenchidas em 2026-08-28 (rascunho v0 -> LaTeX); 2.4-2.5 escritas em 2026-09-17
%% Fonte do rascunho: dissertacao/cap2_rascunho_v0.md
%% Pendencias marcadas com \todo{...}; ver tambem notas no fim do ficheiro.
%% ====================================================================

\chapter{Enquadramento}
\label{cha:enquadramento}

\section{Gestão de energia e o quadro normativo}
\label{sec:normas}

A gestão de energia industrial encontra-se hoje codificada num corpo normativo estruturado em torno da ISO~50001, que especifica os requisitos de um sistema de gestão de energia (SGE), e de um conjunto de normas de suporte, desenvolvidas pelo comité técnico ISO/TC~301, que operacionalizam componentes específicos desse sistema. Para o presente trabalho relevam três documentos: a ISO~50001:2018, adotada em Portugal como NP~EN~ISO~50001:2019~\cite{ISO50001_2018}, que estabelece os requisitos do sistema; a ISO~50006~\cite{ISO50006_2014}, que fornece princípios e orientações para a medição do desempenho energético através de indicadores de desempenho energético (EnPI, na sigla anglófona que a literatura consagrou; IDE na versão portuguesa da norma) e consumos energéticos de referência (EnB, ou CER); e a ISO~50015~\cite{ISO50015_2014}, que estabelece princípios gerais para a medição e verificação (M\&V) do desempenho energético de organizações. A distinção de estatuto entre estes documentos é essencial para a análise que se segue: a ISO~50001 é uma norma de requisitos, certificável, redigida na forma verbal \enquote{deve}; a ISO~50006 e a ISO~50015 são normas de orientação, redigidas na forma \enquote{deverá/poderá}, sem carácter vinculativo --- a própria ISO~50015 declara que os seus princípios não são exigidos pela ISO~50001~\cite[Introdução]{ISO50015_2014}. A ISO~50006 e a ISO~50015 são aqui lidas nas edições de 2014, as consultadas. A ISO~50006 tem desde 2023 uma segunda edição, que substituiu a de 2014 e não foi consultada; o que este capítulo diz sobre o texto da ISO~50006 vale para a edição de 2014 e não se estende, sem verificação, ao quadro atual (nota no fim da Secção~\ref{sec:assimetria-comparacao}).

A ISO~50001 estrutura o SGE sobre o ciclo de melhoria contínua \emph{Plan-Do-Check-Act} (PDCA). Na fase de planeamento, a organização realiza uma avaliação energética (\S6.3), identifica os usos significativos de energia (USE), determina EnPI \enquote{adequados para medir e monitorizar o seu desempenho energético} e que \enquote{permitam demonstrar a melhoria do desempenho energético} (\S6.4), e estabelece EnB a partir da avaliação energética, \enquote{tendo em conta um período de tempo adequado} (\S6.5). Na fase de verificação, a organização monitoriza, mede, analisa e avalia o desempenho (\S9.1): a melhoria do desempenho energético \enquote{deve ser avaliada comparando os valores de IDE [\ldots] com os correspondentes CER}, e a organização \enquote{deve investigar e responder a desvios significativos no desempenho energético} (\S9.1.1). O par EnPI--EnB é, portanto, o mecanismo central através do qual a norma torna o desempenho energético mensurável e a melhoria demonstrável --- e é sobre este par que incide toda a análise desta dissertação.
\alterado{A Tabela~\ref{tab:pdca-50001} localiza as cláusulas no ciclo e
assinala aquelas de que o par depende.}

%% 2026-09-19: o \missingfigure "FIG 2.1" (roda PDCA com as clausulas
%% posicionadas) saia impresso como caixa cinzenta com triangulo vermelho e
%% nao era citado em lado nenhum. Substituido por esta tabela, que faz o
%% mesmo trabalho -- localizar as clausulas no ciclo e separar requisito de
%% orientacao -- e que se pode citar. Versao anterior em
%% _revisao/2026-09-19-cap2-antes.tex.
{\footnotesize
\begin{xltabular}{\linewidth}{@{}>{\raggedright\arraybackslash}p{1.15cm}>{\raggedright\arraybackslash}p{3.5cm}X>{\raggedright\arraybackslash}p{2.35cm}@{}}
\caption[A ISO 50001:2018 pelo ciclo PDCA]{Estrutura da ISO~50001:2018
  pelo ciclo PDCA. A negrito, as quatro cláusulas de que depende o par
  IDE--CER e que esta dissertação examina; entre parênteses, a norma de
  orientação que cada uma convoca e que, por não ser certificável, não
  acrescenta requisito nenhum. A última coluna declara o âmbito: a maior
  parte do sistema de gestão fica fora deste trabalho. Elaboração própria a
  partir de \cite{ISO50001_2018,ISO50006_2014,ISO50015_2014}.}\label{tab:pdca-50001}\\
\toprule
Fase & Cláusula & O que a norma exige & Nesta dissertação \\
\midrule
\endfirsthead
\toprule
Fase & Cláusula & O que a norma exige & Nesta dissertação \\
\midrule
\endhead
\midrule
\multicolumn{4}{r}{\emph{continua na página seguinte}}\\
\endfoot
\bottomrule
\endlastfoot
  \emph{Plan} & 4.1--4.4 Contexto da organização & Determinar o âmbito e as fronteiras do sistema de gestão de energia & --- \\
  \addlinespace
   & 5.1--5.3 Liderança & Política energética; papéis, responsabilidades e autoridades & --- \\
  \addlinespace
   & 6.1--6.2 Riscos, objetivos e metas & Ações face a riscos e oportunidades; objetivos e metas energéticas & --- \\
  \addlinespace
   & \textbf{6.3 Avaliação energética} & Identificar os usos significativos de energia e, para cada um, determinar as variáveis relevantes e o desempenho atual & Secção~\ref{sec:enb-cadeia}, elos 1--3 \\
  \addlinespace
   & \textbf{6.4 Indicadores de desempenho energético} (ISO~50006) & Determinar IDE adequados para medir e monitorizar o desempenho e que permitam demonstrar a melhoria & Secção~\ref{sec:enb-cadeia}, elo 4 \\
  \addlinespace
   & \textbf{6.5 Consumo energético de referência} (ISO~50006) & Estabelecer o CER; normalizar quando as variáveis relevantes afetam significativamente o desempenho; rever o CER em três casos & Secção~\ref{sec:enb-cadeia}, elos 5 e 7 \\
  \addlinespace
   & 6.6 Planeamento da recolha de dados & Especificar que dados recolher, com que frequência e com que qualidade & Capítulo~\ref{cha:metodos} \\
  \midrule
  \emph{Do} & 7.1--7.5 Suporte & Recursos, competência, sensibilização, comunicação, informação documentada & --- \\
  \addlinespace
   & 8.1--8.3 Operação & Planeamento e controlo operacional, conceção, aquisição & --- \\
  \midrule
  \emph{Check} & \textbf{9.1 Monitorização, medição, análise e avaliação} (ISO~50015) & Comparar os valores dos IDE com os CER correspondentes; investigar e responder a desvios significativos & Secção~\ref{sec:enb-cadeia}, elo 7; Capítulo~\ref{cha:diagnostico} \\
  \addlinespace
   & 9.2--9.3 Auditoria interna e revisão pela gestão & Verificar a conformidade e a adequação do sistema & --- \\
  \midrule
  \emph{Act} & 10.1--10.2 Não conformidades e melhoria contínua & Corrigir desvios do sistema e melhorar continuamente o desempenho energético & --- \\
\end{xltabular}
}

A ISO~50006 desenvolve o par EnPI--EnB. Define o processo de medição do desempenho em cinco passos --- obtenção de informação a partir da avaliação energética, identificação dos EnPI, estabelecimento dos EnB, utilização de ambos, e manutenção e ajustamento~\cite[Figura~2]{ISO50006_2014} --- e \alterado{retoma} os conceitos operativos de que o Capítulo~\ref{cha:diagnostico} fará uso intensivo \alterado{--- que a própria ISO~50001 também define (\S3.4.8 a \S3.4.10), com origem declarada na ISO~50015 ---}: fronteira do EnPI (\S4.2.2), \alterado{variável relevante (\S3.14: fator quantificável que afeta o desempenho energético e muda rotineiramente), fator estático (\S3.17: fator identificado que afeta o desempenho energético e não muda rotineiramente)} e normalização (\S3.13: modificação dos dados energéticos para permitir comparação de desempenho \enquote{em condições equivalentes}). A norma tipifica ainda os EnPI em quatro classes --- valor medido, rácio de valores medidos, modelo estatístico e modelo de engenharia (\S4.3.3, Tabela~2) --- tipologia a que regressamos na Secção~\ref{sec:enb} por ser o ponto onde o método de \emph{baseline} usado na refinaria se localiza dentro do vocabulário da própria norma.

A ISO~50015, por fim, enquadra a M\&V como processo: plano de M\&V documentado, fronteiras, seleção de métricas, plano de recolha de dados, estabelecimento e ajustamento de EnB, análise, reporte e incerteza (\S5--\S8). Os seus princípios --- exatidão apropriada e gestão da incerteza, transparência e reprodutibilidade, \alterado{gestão de dados e planeamento da medição,} competência, imparcialidade, \alterado{confidencialidade} e métodos apropriados (\S4) --- são, nas palavras da norma, princípios e não requisitos (\S4.1). No plano internacional coexistem com a família ISO dois protocolos de M\&V com génese no setor dos edifícios e dos contratos de desempenho energético: o IPMVP~\cite{IPMVP_2022} e a ASHRAE Guideline~14~\cite{ASHRAE14_2014}, ambos citados pela bibliografia da própria ISO~50015. A comparação entre estes protocolos e as normas ISO, desenvolvida na Secção~\ref{sec:assimetria}, é instrutiva precisamente porque revela que o silêncio da família ISO quanto a critérios quantitativos de validação é uma escolha de desenho normativo, não uma impossibilidade técnica.

\section{\emph{Baselines} energéticas: conceitos e construção}
\label{sec:enb}

\subsection{O par EnPI--EnB}
\label{sec:enb-par}

O conceito central é simples: um EnPI é uma \enquote{medida ou unidade do desempenho energético, conforme definido pela organização} (ISO~50001, \S3.4.4); o EnB é \enquote{referência(s) quantitativa(s) que fornecem uma referência para comparação do desempenho energético} (\S3.4.7), construída com dados \enquote{obtidos de um período específico de tempo e/ou condições, conforme definido pela organização} (\S3.4.7, Nota~1). A melhoria do desempenho energético demonstra-se comparando o valor do EnPI no período de reporte com o valor de referência do período de \emph{baseline}~\cite[Figuras~1 e~3]{ISO50006_2014}. Quando variáveis relevantes afetam significativamente o desempenho, a comparação direta deixa de ser válida e a organização \enquote{deve realizar a normalização do(s) valor(es) do IDE e do correspondente CER} (ISO~50001, \S6.5) --- tipicamente através de um modelo que descreve a relação entre consumo e variáveis relevantes no período de \emph{baseline}~\cite[\S4.5.1 e Anexo~D]{ISO50006_2014}.

Repare-se, desde já, em duas expressões que a redação normativa usa com naturalidade e que carregam todo o peso estatístico do edifício: a comparação deve fazer-se \emph{em condições equivalentes} (ISO~50006, \S3.13), e o período de \emph{baseline} deve \emph{capturar a variabilidade dos padrões operacionais} (\S4.4.2). Ambas pressupõem que as condições operacionais são descritíveis pelas variáveis escolhidas e que a variabilidade observada num período --- tipicamente doze meses (\S4.4.2) --- é representativa da variabilidade futura. São pressupostos de estacionariedade: não da série de consumo, mas da \emph{relação} entre consumo e variáveis relevantes. Não são enunciados como pressupostos em parte alguma do texto normativo.

\subsection{A cadeia de medição implícita}
\label{sec:enb-cadeia}

Lidas em conjunto, a 50001 (requisitos), a 50006 (construção de EnPI/EnB) e a 50015 (M\&V) desenham uma cadeia de medição: uma sequência de operações que transforma dados de processo em declarações sobre desempenho energético e, no limite, em alarmes e decisões. A cadeia nunca é apresentada como tal em nenhuma das normas --- cada documento descreve o seu troço --- mas a sua reconstrução explícita (Figura~\ref{fig:cadeia-normativa}) é o instrumento analítico deste capítulo, porque permite perguntar, elo a elo, três coisas: o que é exigido, o que é deixado à organização, e que pressuposto tácito o elo precisa de importar para funcionar.

\begin{figure}[htbp]
  \centering
  \begin{tikzpicture}[
      font=\small,
      elo/.style={draw, rounded corners=2pt, align=center, minimum height=11mm, inner sep=2pt, text width=26mm},
      press/.style={align=center, text width=27mm, font=\scriptsize\itshape, text=orange!60!black},
      seta/.style={-{Stealth[length=2.5mm]}, thick},
      node distance=14mm and 4mm
    ]
    % linha superior
    \node[elo] (n1) {1. Avaliação energética\\[1pt] {\scriptsize 50001 \S6.3}};
    \node[elo, right=of n1] (n2) {2. USE e fronteiras\\[1pt] {\scriptsize 50001 \S6.3 · 50006 \S4.2.2}};
    \node[elo, right=of n2] (n3) {3. Variáveis relevantes, fatores estáticos\\[1pt] {\scriptsize 50006 \S4.2.4--4.2.5}};
    \node[elo, right=of n3] (n4) {4. EnPI: tipo e forma\\[1pt] {\scriptsize 50001 \S6.4 · 50006 \S4.3}};
    % linha inferior (fluxo da direita para a esquerda)
    \node[elo, below=of n4] (n5) {5. EnB: período e modelo\\[1pt] {\scriptsize 50001 \S6.5 · 50006 \S4.4}};
    \node[elo, left=of n5] (n6) {6. Comparação normalizada\\[1pt] {\scriptsize 50006 \S4.5, Anexo~D}};
    \node[elo, left=of n6] (n7) {7. Desvio significativo $\to$ ação\\[1pt] {\scriptsize 50001 \S9.1, \S10}};
    % pressupostos tácitos
    \node[press, below=1mm of n3.south] {suficiência das variáveis};
    \node[press, below=1mm of n4.south] {forma funcional adequada};
    \node[press, below=1mm of n5.south] {representatividade do período};
    \node[press, below=1mm of n6.south] {transportabilidade (estacionariedade estrutural)};
    \node[press, below=1mm of n7.south] {calibração e detetabilidade};
    % setas
    \draw[seta] (n1) -- (n2);
    \draw[seta] (n2) -- (n3);
    \draw[seta] (n3) -- (n4);
    \draw[seta] (n4) -- (n5);
    \draw[seta] (n5) -- (n6);
    \draw[seta] (n6) -- (n7);
    % retroação: revisão do EnB
    \draw[seta, dashed] (n7.north) to[out=60, in=120]
      node[midway, above, font=\scriptsize, align=center] {revisão do EnB\\ 50001 \S6.5 · 50006 \S4.6} (n5.north);
  \end{tikzpicture}
  \caption[A cadeia de medição normativa e os seus pressupostos tácitos]{A cadeia de medição implícita na família ISO~50001 e, em itálico, o pressuposto tácito que cada elo importa. Reconstrução própria a partir de \cite{ISO50001_2018,ISO50006_2014,ISO50015_2014}.}
  \label{fig:cadeia-normativa}
\end{figure}

Percorramos os elos onde o texto normativo é mais revelador.

\paragraph{Variáveis relevantes (elo 3).} A ISO~50006 dedica ao problema orientação prática: recomenda gráficos de tendência e diagramas X--Y entre consumo e variável candidata, e observa que \enquote{em muitos casos, uma relação linear simples é adequada para determinar a relevância} (\S4.2.4). O critério de significância, porém, \enquote{é determinado pela organização} (ISO~50001, \S3.4.9, Nota~1) --- a norma diz \emph{como olhar}, mas não fixa nenhum limiar de decisão, nenhum procedimento de teste e, sobretudo, nenhum teste de \emph{suficiência}: nada obriga a organização a demonstrar que o conjunto de variáveis escolhido esgota a variação sistemática do consumo. Nenhuma cláusula distingue a escolha de uma única variável --- a carga processada, no caso em estudo --- da escolha de cinco.

\paragraph{Tipo de EnPI (elo 4).} A Tabela~2 da ISO~50006 é o documento mais interessante de toda a família para o argumento desta dissertação, porque a própria norma enumera, na coluna de observações, as fragilidades de cada tipo. Para o tipo \enquote{valor medido}: não tem em conta os efeitos das variáveis relevantes, \enquote{dando resultados enganadores na maioria das aplicações}. Para o rácio: não considera a carga de base nem efeitos não lineares, sendo \enquote{enganador em instalações com carga de base elevada}. Para o modelo estatístico: \enquote{pode ser impreciso se não for confirmado por testes estatísticos} e \enquote{os modelos devem ser mantidos para assegurar resultados válidos} (\S4.3.3, Tabela~2). \alterado{E, na mesma linha dessa tabela, uma observação que convém ler devagar: \enquote{pode não ser claro se algum erro residual se deve a erro de modelação ou a falta de controlo sobre o consumo de energia}. É o problema desta dissertação enunciado pela própria norma --- o resíduo de uma \emph{baseline} não separa o que o modelo não sabe do que a unidade fez mal --- e enunciado sem que lhe corresponda procedimento nenhum.}\footnote{Traduções do autor a partir do original inglês; as passagens críticas são reproduzidas no original em nota nas ocorrências seguintes.} A norma conhece, portanto, os modos de falha --- enuncia-os --- mas não converte nenhum deles em requisito verificável: não diz \emph{que} testes, \emph{com que limiares}, \emph{com que frequência}, nem \emph{o que fazer} quando falham. A Tabela~\ref{tab:tipologia-enpi} condensa esta tipologia e localiza o método em produção na refinaria --- regressão linear univariada por vetor energético, $E = a \cdot Q + b$ --- na classe \enquote{modelo estatístico}, herdando por escrito os \emph{caveats} que a própria norma lhe associa.

\begin{table}[htbp]
  \centering\small
  \caption[Tipologia de EnPI da ISO 50006 e localização do método em produção]{Tipologia de EnPI da ISO~50006 (\S4.3.3, Tabela~2, condensada; traduções do autor) e localização do método de \emph{baseline} em produção na unidade em estudo.}
  \label{tab:tipologia-enpi}
  \begin{tabularx}{\textwidth}{@{}l X X@{}}
    \toprule
    \textbf{Tipo de EnPI} & \textbf{Aplicação típica (norma)} & \textbf{Fragilidades enunciadas pela própria norma} \\
    \midrule
    Valor medido & Reduções absolutas; monitorização de custos & Ignora variáveis relevantes; \enquote{resultados enganadores na maioria das aplicações} \\
    Rácio de valores medidos & Eficiência com uma só variável e carga de base pequena; \emph{benchmarking} & Ignora carga de base e efeitos não lineares \\
    \textbf{Modelo estatístico} $\leftarrow$ \emph{método em produção} & Várias variáveis relevantes; carga de base; normalização & Impreciso se não confirmado por testes estatísticos; \alterado{pode não ser claro se o erro residual vem do modelo ou da falta de controlo do consumo;} exige manutenção para assegurar validade \\
    Modelo de engenharia & Variáveis numerosas ou interdependentes; fase de projeto & Exige manutenção para assegurar validade \\
    \bottomrule
  \end{tabularx}
\end{table}

\paragraph{Estabelecimento e teste do EnB (elo 5).} A cláusula \S4.4.3 da ISO~50006 é o ponto de máxima aproximação de toda a família a um requisito de validação --- e, por isso mesmo, o ponto onde o silêncio é mais nítido: \enquote{Se apropriado, o EnB deverá ser testado quanto à validade, para assegurar que constitui uma referência apropriada para comparação. Quando se usam modelos, a validade do EnB pode ser determinada usando testes estatísticos como o valor-$p$, o teste~$F$ ou o coeficiente de determinação [\ldots]. Se se determinar que um modelo não é válido, a organização deverá considerar ajustar o EnB ou determinar um novo modelo}.\footnote{No original: \enquote{If appropriate, the EnB should be tested for validity to ensure that it is an appropriate reference for comparison. When models are used, the validity of the EnB can be determined using statistical tests such as the P-Value, F-Test or the coefficient of determination [\ldots].} \cite[\S4.4.3]{ISO50006_2014}.} Três observações: \begin{enumerate*}[(i)] \item \enquote{se apropriado} e \enquote{deverá} --- recomendação condicional, não requisito; \item os testes são nomeados, mas nenhum limiar de aceitação é fixado --- \alterado{e o Anexo~D, que é onde a norma mais perto chega de uma exigência, diz que cada variável explicativa \enquote{terá de cumprir um certo valor-$p$} (\S D.2) sem dizer qual}; \item as três estatísticas descrevem o ajuste no período de treino: o $R^2$ é uma medida descritiva, que se calcula sem pressupostos sobre os resíduos, e a validade de um valor-$p$ ou de um teste~$F$ depende do teste usado e das suas condições (na forma clássica, independência, homoscedasticidade e normalidade dos resíduos), que a norma não manda verificar\end{enumerate*}. Mesmo quando válidas, nenhuma delas diz se a referência prevê o período seguinte, se as bandas têm a cobertura que se lhes atribui ou se um desvio seria detetado. A norma sugere, como teste de validade, estatísticas de ajuste que não medem aquilo de que depende o uso da referência.

\paragraph{Revisão do EnB (elo 7, retroação).} A ISO~50001 fixa três gatilhos de revisão do EnB: quando os EnPI \enquote{deixaram de refletir o desempenho energético da organização}; quando \enquote{ocorram mudanças significativas nos fatores estáticos}; ou \enquote{de acordo com um método predeterminado} (\S6.5). A ISO~50006 acrescenta os ajustamentos não-rotineiros e nota que os métodos predeterminados \enquote{tipicamente repõem o EnB em intervalos definidos} (\S3.1, Nota~3), sancionando explicitamente a prática da recalibração periódica. O que nenhum dos gatilhos fornece é um \emph{mecanismo de deteção}: o primeiro é circular (como sabe a organização que os EnPI deixaram de refletir o desempenho, se o instrumento de medição do desempenho são os próprios EnPI?); o segundo pressupõe que a mudança de fator estático é conhecida e identificada; o terceiro --- o calendário --- não é sensível a mudança nenhuma. O Capítulo~\ref{cha:diagnostico} mostrará a consequência empírica: a recalibração anual, que a norma admite como método predeterminado, absorve no intercepto a deriva estrutural que devia alarmar.

\subsection{Onde as normas param e a implementação começa}
\label{sec:enb-camadas}

A análise anterior obriga a uma distinção de três camadas, que atravessa toda a dissertação e sem a qual a crítica seria imprecisa. Primeira camada, \textbf{o que as normas exigem} (ISO~50001, \enquote{deve}): que existam EnPI e EnB, que a melhoria se demonstre por comparação entre ambos, que se normalize quando as variáveis relevantes afetam significativamente o desempenho, que os desvios significativos sejam investigados, e que o EnB seja revisto nos três gatilhos. Segunda camada, \textbf{o que as normas sugerem} (ISO~50006/50015, \enquote{deverá}, anexos informativos): a tipologia de EnPI, testes de validade nomeados sem limiares, o período de \emph{baseline} de doze meses, a gestão da incerteza \enquote{ao nível apropriado} (ISO~50015, \S4.2) e os diagnósticos possíveis da incerteza de modelo --- \enquote{estatísticas~$t$, valores de $R^2$, valores-$p$, níveis de confiança, limites de predição do modelo} (ISO~50015, \S7, Nota) --- \enquote{na medida do praticável} (\S7). Terceira camada, \textbf{o que a implementação escolhe} --- e que nenhuma norma prescreve nem proíbe: no caso em estudo, a forma $E = a \cdot Q + b$ com uma única variável explicativa; a estimação por mínimos quadrados ordinários; as bandas de alarme em $\pm 2\sigma$, com $\sigma$ estimado a partir dos resíduos de treino; e a recalibração anual por calendário.

Esta atribuição precisa é essencial: o objeto da crítica do Capítulo~\ref{cha:diagnostico} não é a conformidade. A forma do modelo e a revisão periódica não são, por si, excluídas por nenhuma das cláusulas analisadas. Três coisas ficam separadas: as normas não prescrevem nenhum teste quantitativo de adequação; a demonstração de que os EnPI e as EnB são adequados cabe à organização (ISO~50001, \S6.4 e \S6.5); e saber se uma auditoria aceitaria um método sem essa demonstração é uma questão que esta dissertação não avalia. O objeto da crítica é o facto de o espaço permitido conter, sem qualquer barreira normativa, implementações cujos pressupostos estatísticos são sistematicamente violados pelos dados a que se aplicam --- e de as normas não fornecerem ao operador nenhum instrumento que o alerte para isso. Dito de outro modo: a ausência de uma barreira normativa não é evidência de validade estatística, e as normas analisadas não contêm o mecanismo que permitiria distinguir uma da outra. \alterado{A Figura~\ref{fig:enpi-enb-meta} mostra onde passa essa fronteira.}

%% 2026-09-19: o \missingfigure "FIG 2.3" foi substituido por um desenho
%% proprio em TikZ, sem dados e sem reproducao de arte da ISO.
%% Figura: o conceito EnPI/EnB/meta e a fronteira entre o que a norma
%% desenha e o que a implementacao acrescenta. Desenho proprio (2026-09-19);
%% substitui o \missingfigure "FIG 2.3". Nao reproduz arte da ISO: e um
%% esquema conceptual sem dados, construido a partir do texto das clausulas.
\begin{figure}[htbp]
  \centering
  \begin{tikzpicture}[x=1.2cm, y=1cm]
    % ---- banda de alarme (extensao da implementacao) ----
    \begin{scope}[on background layer]
      \fill[corFundo] (0.15,3.0) rectangle (10.9,3.8);
    \end{scope}
    \draw[corCinzaC, line width=.5pt, dash pattern=on 2pt off 1.4pt]
      (0.15,3.8) -- (10.9,3.8);
    \draw[corCinzaC, line width=.5pt, dash pattern=on 2pt off 1.4pt]
      (0.15,3.0) -- (10.9,3.0);
    \draw[corCinzaC, line width=.5pt] (11.0,3.0) -- (11.15,3.0)
      -- (11.15,3.8) -- (11.0,3.8);
    \node[ntmini, anchor=west] at (11.2,3.4) {$\pm 2\sigma$};

    % ---- eixos ----
    \draw[corCinza, line width=.6pt] (0.15,1.75) -- (11.0,1.75);
    \draw[corCinza, line width=.6pt] (0.15,1.75) -- (0.15,4.95);
    \node[ntrot, rotate=90, anchor=south] at (-0.12,3.35)
      {EnPI (consumo normalizado)};
    \node[ntmini, anchor=north east] at (11.0,1.67) {tempo};

    % ---- fronteira entre periodos ----
    \draw[corTinta, line width=.6pt, dash pattern=on 2pt off 1.6pt]
      (5.5,1.75) -- (5.5,4.80);
    \node[ntrot, anchor=south] at (2.80,4.92) {período de referência (EnB)};
    \node[ntrot, anchor=south] at (8.25,4.92) {período de reporte};

    % ---- modelo (EnB) ----
    \draw[corFG, line width=1pt] (0.15,3.4) -- (10.9,3.4);
    \node[ntmini, text=corFG, anchor=north east] at (10.85,3.34)
      {EnB: modelo estimado à esquerda};

    % ---- meta ----
    \draw[corCinza, line width=.6pt, dash pattern=on 3pt off 2pt]
      (0.15,2.40) -- (10.9,2.40);
    \node[ntmini, anchor=south west] at (0.28,2.43) {meta energética};

    % ---- consumo observado ----
    \draw[corTinta, line width=.7pt]
      (0.20,3.34) -- (0.90,3.56) -- (1.60,3.20) -- (2.30,3.52) --
      (3.00,3.28) -- (3.70,3.56) -- (4.40,3.24) -- (5.10,3.46) -- (5.50,3.38)
      -- (6.20,3.52) -- (6.90,3.44) -- (7.60,3.66) -- (8.30,3.62) --
      (9.00,3.80) -- (9.70,3.90) -- (10.40,4.08) -- (10.90,4.15);
    \node[ntmini, anchor=south east] at (10.9,4.22) {valor do EnPI};

    % ---- chamadas ----
    \draw[corCinza, line width=.4pt] (1.25,4.28) -- (1.25,3.84);
    \node[circle, draw=corTinta, line width=.4pt, fill=white, inner sep=1pt,
          font=\tiny\bfseries, text=corTinta] at (1.25,4.38) {1};
    \node[circle, draw=corTinta, line width=.4pt, fill=white, inner sep=1pt,
          font=\tiny\bfseries, text=corTinta] at (5.5,4.62) {2};

    % ---- notas ----
    \node[ntproposta, text width=5.6cm, anchor=north west, align=left]
      at (0.15,1.35)
      {\textbf{1.} As bandas de alarme não constam de nenhuma figura
       normativa: são extensão da implementação.};
    \node[ntproposta, text width=5.6cm, anchor=north west, align=left]
      at (5.35,1.35)
      {\textbf{2.} Tudo o que está à direita pressupõe que a relação
       estimada à esquerda continua válida.};
  \end{tikzpicture}
  \caption[O conceito EnPI/EnB/meta e a extensão de implementação]{O conceito
   EnPI/EnB/meta e a fronteira entre o que a norma desenha e o que a
   implementação acrescenta. A comparação entre o valor do indicador e a
   referência é o mecanismo que a ISO~50001 exige (\S9.1.1); a banda de
   alarme em torno da referência não aparece em nenhuma figura normativa e é
   escolha da organização. O Anexo~D da ISO~50006 usa o modelo estimado no
   período de referência para prever o período de reporte e lê a diferença
   como melhoria de desempenho, sem enunciar a transportabilidade da relação
   como pressuposto. Esquema próprio, sem dados, a partir de
   \cite{ISO50001_2018,ISO50006_2014}.}
  \label{fig:enpi-enb-meta}
\end{figure}


\section{A assimetria normativa}
\label{sec:assimetria}
% [REV-2026-09-13][C2-05][P2][ESTRUTURA]
\revisaotese{C2-05}{P2}{ESTRUTURA}
{Fundir passagens repetidas e separar texto normativo de interpretação}
{As Secs. 2.2.2, 2.2.3 e 2.3.1 repetem a mesma crítica. Manter uma
explicação curta e uma única matriz requisito/orientação/escolha local,
com edição e cláusula. Assinalar "interpretação desta dissertação" para
pressupostos reconstruídos: a norma não declara literalmente todos os
elos do diagrama. Nas citações, preferir paráfrases com cláusula e
reservar transcrições breves para a distinção verbal realmente analisada.}%
% [/REV-2026-09-13]

\subsection{Formulação}
\label{sec:assimetria-formulacao}

Chamamos assimetria normativa ao desequilíbrio, documentado cláusula a cláusula na secção anterior, entre a densidade prescritiva das normas quanto à \emph{existência} dos instrumentos de medição do desempenho energético e o seu silêncio quanto à \emph{validade} desses instrumentos. A família ISO~50001 obriga a que haja EnPI, EnB, normalização, investigação de desvios e revisão da referência; nomeia, inclusive, os instrumentos estatísticos com que a validade poderia ser aferida; mas em nenhum ponto converte a aferição em requisito, fixa limiares de aceitação, exige a verificação dos pressupostos dos modelos ou fornece critérios operacionais para detetar que uma referência deixou de ser válida. O resultado prático é que nenhuma das cláusulas analisadas fixa um teste que impeça um sistema de gestão certificado de operar sobre \emph{baselines} cujos pressupostos os dados violam: a adequação fica por demonstrar pela organização e por julgar pela auditoria, o que aqui não se avalia. O Capítulo~\ref{cha:diagnostico} examina essa situação nos vetores energéticos da unidade em estudo (Secção~\ref{sec:caso-estudo}).

A assimetria tem uma segunda face, menos evidente: as normas embutem pressupostos estatísticos que nunca enunciam. A comparação \enquote{em condições equivalentes} pressupõe a suficiência das variáveis relevantes; a extrapolação do modelo de \emph{baseline} para o período de reporte pressupõe a estabilidade estrutural da relação; o período de referência de doze meses pressupõe que um ciclo anual é representativo; os gatilhos de revisão pressupõem que a mudança é observável por quem opera. \alterado{O Anexo~D da ISO~50006 é o lugar onde o pressuposto da transportabilidade mais trabalha e menos se declara: o modelo estimado no período de referência é usado para prever o período de reporte, e a diferença entre o previsto e o observado \enquote{indicará se ocorreu uma melhoria de desempenho energético} (\S D.1). O anexo ressalva que os valores previstos e os observados não coincidem exatamente; não ressalva que a relação estimada à esquerda possa ter deixado de valer à direita.} Nenhum destes pressupostos é neutro num processo contínuo sujeito a \emph{fouling}, desativação de catalisador e deslocamento gradual do espaço operacional --- a especificidade desenvolvida na Secção~\ref{sec:processos-continuos}.

\subsection{Matriz de rastreabilidade: cláusula, silêncio, pressuposto, teste}
\label{sec:assimetria-matriz}

A Tabela~\ref{tab:matriz-rastreabilidade} formaliza a análise: para cada cláusula relevante, o que é exigido ou sugerido, o que fica por especificar, o pressuposto tácito que a cláusula importa, e o \emph{endpoint} do protocolo de diagnóstico (Secção~\ref{sec:protocolo-baselines}) que o operacionaliza. É esta última coluna que dá ao Capítulo~\ref{cha:diagnostico} rastreabilidade completa norma\,$\to$\,resultado.

\begin{table}[htbp]
  \centering\footnotesize
  \caption[Matriz de rastreabilidade cláusula--silêncio--pressuposto--teste]{Matriz de rastreabilidade: cláusula normativa, silêncio, pressuposto tácito e \emph{endpoint} de diagnóstico que o testa (E1--E6 e bateria descritiva D; ver Secção~\ref{sec:protocolo-baselines}).}
  \label{tab:matriz-rastreabilidade}
  \begin{tabularx}{\textwidth}{@{}c >{\raggedright\arraybackslash}p{2.1cm} X X X c@{}}
    \toprule
    \# & \textbf{Cláusula} & \textbf{Exige/sugere} & \textbf{Deixa por especificar} & \textbf{Pressuposto tácito} & \textbf{Teste} \\
    \midrule
    1 & 50001 \S6.4; 50006 \S4.3 & EnPI \enquote{adequados para medir e monitorizar} & Critério de adequação; teste de conteúdo preditivo & O EnPI mede desempenho, não ruído & E2 \\
    2 & 50001 \S3.4.9~N1; 50006 \S4.2.4 & Determinar variáveis relevantes (métodos gráficos) & Limiar de significância; teste de \emph{suficiência} do conjunto & As variáveis descrevem o estado operacional & E6 \\
    3 & 50006 \S4.4.2 & Período que \enquote{capture a variabilidade} (tip.\ 12 meses) & Verificação de representatividade & Um ciclo anual antecipa o futuro & E3, E4 \\
    4 & 50006 \S4.4.3 & \enquote{Testar a validade} do EnB ($p$, $F$, $R^2$ nomeados) & \enquote{Se apropriado}; sem limiares; sem verificação das condições dos testes & O ajuste no treino indica desempenho fora dele; condições da inferência clássica cumpridas & E1, E2, D \\
    5 & 50001 \S6.5; 50006 \S4.5, An.~D & Normalizar por variáveis relevantes via modelo & Validação do modelo de normalização; incerteza & A relação transporta-se do treino para o reporte & E3 \\
    6 & 50001 \S9.1.1 & Investigar \enquote{desvios significativos} & Definição de significativo (remetida para a organização) & Existe calibração conhecida do desvio & E4 \\
    7 & 50001 \S6.5~a)--c); 50006 \S4.6, \S3.1~N3 & Rever o EnB (3 gatilhos; \emph{reset} periódico sancionado) & Mecanismo de \emph{deteção} de invalidade; gatilho~a) circular & Mudanças invalidantes são conhecidas ou visíveis & E5 \\
    8 & 50015 \S7 + Nota & Identificar e quantificar incerteza ($t$, $R^2$, $p$, limites de predição nomeados) & Obrigatoriedade; limiares; reporte junto ao resultado & A incerteza declarada é gerível caso a caso & E4 \\
    9 & 50015 \S4.1, Introd. & Princípios de M\&V (exatidão, reprodutibilidade) & \enquote{Princípios, não requisitos}; não exigidos pela 50001 & A boa prática é adotada voluntariamente & --- \\
    10 & 50006 \S4.2.6.4 & Rever qualidade dos dados; examinar \emph{outliers} & Procedimento e critério de exclusão; efeito sobre $\sigma$ & \emph{Outliers} são raros e identificáveis & D \\
    \bottomrule
  \end{tabularx}
\end{table}

A leitura vertical da coluna dos pressupostos produz a lista compacta que estrutura o resto da dissertação: \textbf{suficiência} (as variáveis descrevem o estado), \textbf{forma} (a classe de modelo captura a relação), \textbf{representatividade} (o período de treino antecipa o futuro), \textbf{transportabilidade} (a relação não muda entre treino e aplicação), \textbf{calibração} (os desvios têm escala conhecida) e \textbf{detetabilidade} (a invalidação é observável). Os resultados do Capítulo~\ref{cha:diagnostico} organizam-se por estes pressupostos, com uma ressalva: nem todos os testaram com o mesmo alcance, e a detetabilidade real ficou por estimar.

\subsection{O silêncio como escolha: comparação com o IPMVP e a ASHRAE Guideline 14}
\label{sec:assimetria-comparacao}

Poderia argumentar-se que o silêncio quantitativo das normas ISO é inevitável --- que critérios de validação universais não podem ser fixados para organizações de todos os tipos e dimensões; é, de resto, a justificação implícita do \enquote{conforme definido pela organização} que percorre os três documentos. A comparação com os protocolos de M\&V do domínio dos edifícios desmente a inevitabilidade, e fá-lo em três frentes distintas.

Primeira frente: \emph{limiares de calibração de modelos}. A ASHRAE Guideline~14 exige, como requisito de conformidade da via de simulação calibrada, um erro médio normalizado (NMBE) máximo de $\pm 5\,\%$ e um coeficiente de variação do erro quadrático médio, CV(RMSE), máximo de $15\,\%$ face a dados de calibração mensais --- $\pm 10\,\%$ e $30\,\%$, respetivamente, para dados horários~\cite[\S4.3.2.4\,f), \S5.3.3.3.10]{ASHRAE14_2014}. As orientações de M\&V do programa federal norte-americano (FEMP) adotam exatamente os mesmos limiares, por remissão declarada para a ASHRAE~\cite[\S4.5.3.2, Tabela~4-2]{FEMP_2015}. Segundo \textcite{RuizBandera2017}, o IPMVP, que desde as edições de 2009 remete os detalhes técnicos de calibração para a Guideline~14, recomenda também, tal como o \emph{ASHRAE Handbook}, que o $R^2$ de um modelo calibrado não seja inferior a $0{,}75$. A recomendação é relatada por essa fonte secundária e não foi confirmada no texto do IPMVP; não se apresenta como critério atual verificado.

Segunda frente, mais próxima do caso em estudo: \emph{limiares sobre o próprio modelo de regressão de baseline}. A via prescritiva de edifício inteiro da Guideline~14 --- precisamente o cenário de uma \emph{baseline} de regressão sobre dados de contagem, análogo estrutural do método em produção na refinaria --- impõe que o modelo de \emph{baseline} tenha CV(RMSE) máximo de $20\,\%$ para energia (escalando para $25\,\%$ e $30\,\%$ com períodos de reporte mais longos), que o período de \emph{baseline} cubra pelo menos doze meses contínuos com um mínimo de nove pontos válidos, e que \emph{nenhum ponto seja eliminado} do período de \emph{baseline}~\cite[\S4.3.2.1]{ASHRAE14_2014}; a via de simulação calibrada admite exclusões, mas limita-as a $25\,\%$ dos dados medidos e exige a documentação das razões~\cite[\S4.3.2.4]{ASHRAE14_2014}. O FEMP, na Opção~C (análise de regressão sobre dados de faturação --- de novo o análogo do caso Galp), formula-o sem ambiguidade: \enquote{a validade estatística dos modelos de regressão finais tem de ser demonstrada}~\cite[\S4.4]{FEMP_2015}.\footnote{No original: \enquote{Statistical validity of the final regression models must be demonstrated.} \cite[\S4.4]{FEMP_2015}.}

Terceira frente: \emph{incerteza com critério de aceitação}. A Guideline~14 não se limita a pedir que a incerteza seja reportada com cada relatório de poupanças: fixa um teto --- a incerteza tem de ser inferior a $50\,\%$ das poupanças anuais reportadas, a um nível de confiança de $68\,\%$~\cite[\S4.3.2.4\,h)]{ASHRAE14_2014}.

Independentemente do juízo sobre estes limiares em concreto --- a literatura discute a sua suficiência estatística, e o Capítulo~\ref{cha:metodos} adotará critérios de outra natureza, sobre desempenho fora do treino, calibração e deteção, que não se comparam com estes limiares de ajuste --- a sua simples existência demonstra que protocolos com ambição de aplicação geral \emph{conseguem} fixar critérios quantitativos de aceitação de modelos, regras de tratamento de dados e tetos de incerteza. A família ISO~50001 escolheu não o fazer. A Tabela~\ref{tab:comparacao-frameworks} sintetiza a comparação.

\begin{table}[htbp]
  \centering\footnotesize
  \caption[Comparação de quadros normativos quanto a requisitos de validação]{Comparação de quadros normativos e protocolos de M\&V quanto a requisitos de validação de modelos de \emph{baseline}. ASHRAE Guideline~14 na edição de 2014; o conteúdo da coluna IPMVP é relatado por \textcite{RuizBandera2017}, sem consulta da edição corrente do protocolo.}
  \label{tab:comparacao-frameworks}
  \begin{tabularx}{\textwidth}{@{}>{\raggedright\arraybackslash}p{3.2cm} X X X X@{}}
    \toprule
    \textbf{Dimensão} & \textbf{ISO 50001:2018} & \textbf{ISO 50006:2014 / 50015:2014} & \textbf{IPMVP} & \textbf{ASHRAE G14} \\
    \midrule
    Estatuto & Requisitos (certificável) & Orientação & Protocolo voluntário & \emph{Guideline} \\
    Validação do modelo de \emph{baseline} & --- & \enquote{deverá, se apropriado} (50006 \S4.4.3) & Sim, por adesão ao protocolo & Sim (requisito de conformidade) \\
    Limiares quantitativos & Nenhum & Nenhum (testes/diagnósticos apenas nomeados) & $R^2 \geq 0{,}75$ (recomendação relatada, simulação calibrada); detalhe remetido para a G14 & NMBE $\pm5/\pm10\,\%$, CV(RMSE) $15/30\,\%$ (calibração, \S4.3.2.4); CV(RMSE) $\leq 20$--$30\,\%$ no modelo de regressão de base (\S4.3.2.1) \\
    Tratamento de exclusões e \emph{outliers} & --- & Rever, sem critério (50006 \S4.2.6.4) & Justificação exigida & Proibidas na via prescritiva; $\leq 25\,\%$, documentadas, na simulação \\
    Verificação dos pressupostos dos testes & --- & --- & Parcial & Parcial \\
    Incerteza reportada com o resultado & --- & \enquote{na medida do praticável} (50015 \S7) & Sim (Apêndice~B) & Sim, com teto: $<50\,\%$ das poupanças a $68\,\%$ de confiança \\
    Deteção de invalidação da \emph{baseline} & 3 gatilhos qualitativos (\S6.5) & Qualitativa (50006 \S4.6; 50015 \S5.10.2) & Ajustes não-rotineiros & --- \\
    \bottomrule
  \end{tabularx}
\end{table}

Duas notas de rigor. Primeira: os limiares ASHRAE/IPMVP/FEMP foram desenvolvidos para modelos de energia de edifícios, com sazonalidade meteorológica dominante; a sua transposição direta para processo contínuo não é defendida aqui --- o que se importa é o \emph{princípio} de critérios de aceitação explícitos, não os números. Segunda: mesmo estes protocolos são imperfeitos --- \textcite{RuizBandera2017} documentam inconsistências de definição entre eles e entre as suas sucessivas edições (a confusão sistemática entre MBE e NMBE, fórmulas divergentes, limiares que mudam de versão para versão), e todos são omissos quanto à verificação dos pressupostos distribucionais dos seus próprios índices e quanto à não-estacionariedade estrutural. O \emph{gap} identificado no Capítulo~\ref{cha:revisao} é, portanto, comum a todo o panorama normativo, não exclusivo da família ISO; a família ISO é simplesmente o caso extremo, por não fixar sequer os índices.

\paragraph{Nota sobre as edições analisadas.} A presente análise reporta-se à ISO~50006:2014 e à ISO~50015:2014, as edições a que houve acesso e, presumivelmente, as que enquadram o método de \emph{baseline} em produção na instalação em estudo. Existe uma segunda edição da ISO~50006, de 2023, que substituiu a de 2014, harmonizada com a ISO~50001:2018, que introduz, entre outras alterações, uma cláusula dedicada à incerteza de modelo (\S8.2)~\cite{ISO50006_2023}; não foi possível consultá-la no âmbito deste trabalho. Regista-se como limitação (Secção~\ref{sec:limitacoes}): as afirmações desta secção sobre o silêncio normativo estão verificadas contra as edições de 2014, e a possibilidade de a edição de 2023 mitigar parcialmente algum dos silêncios identificados não pôde ser excluída nem confirmada. Assinale-se, porém, que a ISO~50006:2023 mantém o estatuto de norma de orientação, pelo que o argumento central --- a ausência de \emph{requisitos} de validação no único documento certificável da família, a ISO~50001 --- não depende dessa edição.

\subsection{Síntese}
\label{sec:assimetria-sintese}

O quadro normativo obriga a medir e a demonstrar melhoria, sugere sem obrigar os instrumentos de validação, e é silencioso sobre os pressupostos estatísticos de que a validade desses instrumentos depende. Para um processo estacionário e bem comportado, o silêncio é inconsequente. A questão a que o resto da dissertação responde empiricamente é o que acontece quando o processo não é estacionário --- quando o \emph{fouling}, as transições de consumidores e o deslocamento do espaço operacional fazem parte da física da unidade (Secção~\ref{sec:processos-continuos}) --- e o espaço de implementação permitido pela norma é preenchido pela escolha mais simples que nele cabe: uma regressão univariada estática, com bandas de $\pm 2\sigma$, recalibrada por calendário (Secção~\ref{sec:caso-estudo}).

\section{Especificidade dos processos contínuos}
\label{sec:processos-continuos}

Num edifício, o que desloca a relação entre consumo e variável explicativa
é quase sempre exterior ao sistema e periódico: a temperatura exterior, a
ocupação, o calendário. Num processo contínuo de refinação, a relação move-se
também por dentro. O equipamento degrada-se enquanto opera, a matéria-prima
muda de composição de um navio para o seguinte, a unidade troca calor e
vapor com as vizinhas, e as intervenções que repõem o desempenho fazem-no aos
saltos. Estes mecanismos têm escalas de tempo diferentes, pedem respostas
diferentes e, sobretudo, nem todos são ineficiência. É esta a diferença que
torna os pressupostos da Secção~\ref{sec:assimetria-matriz} verificáveis em
vez de razoáveis.

\paragraph{Degradação com recuperação: a incrustação.} Numa unidade de
destilação atmosférica, o crude é aquecido primeiro numa bateria de
permutadores que recupera calor dos produtos e dos refluxos circulantes da
coluna, e só depois no forno, que fornece o que a bateria não entregou~\cite{Gary2007}.
Os depósitos que o crude forma nas superfícies de permuta reduzem o calor
recuperado; a temperatura à entrada do forno desce e o forno queima mais
combustível para a mesma carga e o mesmo corte~\cite{ColettiMacchietto2011}.
A velocidade de deposição depende da temperatura da parede e da tensão de
corte junto a ela, o que faz com que a mesma bateria incruste a ritmos
diferentes consoante o caudal e o crude~\cite{Wilson2017}. O depósito envelhece e
endurece, e a limpeza, que interrompe a permuta e tem custo, é agendada num
compromisso entre a energia perdida e o custo de parar o
permutador~\cite{Ishiyama2011}. O resultado, visto a partir do consumo
específico do forno, é um dente de serra: subida gradual ao longo de semanas
ou meses, queda na limpeza, recomeço. É um mecanismo físico previsível e, ao
mesmo tempo, uma perda de eficiência real que a operação pode adiar ou
antecipar. Um modelo que só conhece a carga atribui a subida ao resíduo, e um
modelo reajustado todos os anos absorve-a na ordenada na origem.

\paragraph{Mudança da matéria-prima.} A refinação processa crudes diferentes
em sequência, e a composição da carga determina quanto dela vaporiza na zona
de \emph{flash}, que rendimentos saem em cada corte e que calor é preciso
para os obter~\cite{Gary2007}. A mesma tonelada de crude não pede a mesma
energia. Parte deste efeito é observável: a densidade, as condições de
\emph{flash} e o rendimento em resíduo são medidos ou calculados
diariamente, e entram, por exemplo, no índice de intensidade energética da
refinaria. A composição altera também a tendência para incrustar, pelo que
este mecanismo e o anterior aparecem misturados nos dados.

\paragraph{Modo de operação e fronteira.} A mesma unidade pode ser operada para
favorecer um produto ou outro, com refluxos, vapor de \emph{stripping} e
temperaturas de corte diferentes, e em regimes de carga reduzida em que o
consumo deixa de ser proporcional à carga porque uma parte dele não depende
dela. Além disso, uma unidade de destilação não é uma ilha energética: recebe
e cede calor e vapor à rede da refinaria, produz vapor de um nível e consome
de outros. O que se contabiliza como consumo da unidade depende do estado da
rede e da convenção de fronteira, e pode mudar sem que nada mude dentro dela.

\paragraph{Desativação de catalisador.} Nas unidades de conversão com reator,
fora da destilação, o catalisador perde atividade ao longo do ciclo por
deposição de coque, envenenamento e sinterização~\cite{Bartholomew2001}. A
compensação habitual é subir a temperatura de reação, com mais energia
fornecida por tonelada à medida que o ciclo avança, até à regeneração ou
substituição. É o mesmo padrão do dente de serra, com uma escala de tempo mais
longa. Não se aplica à unidade estudada, mas aplica-se a várias das
dezasseis unidades que a plataforma cobre.

\paragraph{Intervenções e medição.} Limpezas, paragens para manutenção,
alterações de equipamento e recalibrações de instrumentos produzem mudanças
em degrau. Umas repõem desempenho, outras alteram a fronteira ou a escala da
medição, e nenhuma é visível para um modelo que não receba o registo do
evento. A distinção importa: uma melhoria real, uma limpeza e um instrumento
substituído produzem resíduos com a mesma forma.

A destilação atmosférica e a de vácuo estão entre as maiores consumidoras de
energia de uma refinaria, e a manutenção dos fornos e a recuperação de calor
nas baterias de pré-aquecimento figuram entre as medidas de poupança
identificadas para estas unidades~\cite{WorrellGalitsky2005}. São também,
pelas razões desta secção, o sítio onde uma referência estática mais
facilmente falha. A Tabela~\ref{tab:mecanismos-pressupostos} fecha a secção: para cada
mecanismo, o efeito esperado sobre a relação consumo--carga, o que dele se
observa, o que se pode confundir com ele e o pressuposto da
Secção~\ref{sec:assimetria-matriz} que viola. A última coluna é a ligação
entre a física desta secção e os testes do Capítulo~\ref{cha:diagnostico}.

\begin{table}[htbp]
  \centering\footnotesize
  \caption[Mecanismos de processo contínuo e pressupostos que violam]{Mecanismos
  de não estacionariedade em processo contínuo, o que deles se observa e o
  pressuposto tácito que violam. CDU: presente na unidade estudada. Os
  pressupostos são os da Secção~\ref{sec:assimetria-matriz}.}
  \label{tab:mecanismos-pressupostos}
  \begin{tabularx}{\textwidth}{@{}>{\raggedright\arraybackslash}p{2.3cm} >{\raggedright\arraybackslash}X >{\raggedright\arraybackslash}X >{\raggedright\arraybackslash}X >{\raggedright\arraybackslash}p{2.1cm}@{}}
    \toprule
    \textbf{Mecanismo} & \textbf{Efeito na relação consumo--carga} & \textbf{O que se observa} & \textbf{Com que se confunde} & \textbf{Pressuposto violado} \\
    \midrule
    Incrustação da bateria de pré-aquecimento (CDU) & Mais combustível no forno para a mesma carga; subida gradual e queda na limpeza & Temperaturas e caudais dos permutadores, a partir dos quais se pode estimar a resistência de incrustação; o estado não é medido & Mudança de crude; excesso de ar no forno; ineficiência operacional & Suficiência; transportabilidade \\
    \addlinespace
    Composição do crude (CDU) & Energia por tonelada diferente de crude para crude; muda também a tendência para incrustar & Densidade, condições de \emph{flash}, rendimento em resíduo; plano de crudes & Incrustação; mudança de modo de operação & Suficiência; forma \\
    \addlinespace
    Modo de operação e carga reduzida (CDU) & Refluxos e vapor de \emph{stripping} diferentes; parte do consumo deixa de acompanhar a carga & Temperaturas de corte, vapor de \emph{stripping}, carga & Dias atípicos; ineficiência & Forma; representatividade \\
    \addlinespace
    Estado da rede de vapor e integração térmica (CDU e refinaria) & O consumo líquido de um nível de vapor muda de sinal sem mudança na unidade & Balanço de utilidades por nível de pressão & Variação real de consumo & Suficiência; calibração \\
    \addlinespace
    Desativação de catalisador (outras unidades) & Mais energia por tonelada ao longo do ciclo, até à regeneração & Temperatura de reação; dias desde a regeneração & Mudança de carga ou de severidade & Transportabilidade; representatividade \\
    \addlinespace
    Limpezas, paragens, recalibrações & Degraus na relação ou na escala da medição & Registos de manutenção e de instrumentação, fora dos dados analisados & Melhoria real do desempenho & Transportabilidade; detetabilidade \\
    \bottomrule
  \end{tabularx}
\end{table}

A Secção~\ref{sec:corpus-b} mostra que a literatura sobre indicadores e
referências nomeia estas causas raramente: relata o efeito no modelo e
quase nunca o mecanismo no processo. É uma razão adicional para as fixar
aqui, antes dos dados.

\section{O caso de estudo: a refinaria de Sines e a unidade de destilação atmosférica}
\label{sec:caso-estudo}
A Figura~\ref{fig:fronteira-cdu} fixa o vocabulário: onde passa a fronteira do indicador, por onde entra cada vetor energético e em que ponto se mede.

\begin{figure}[htbp]
  \centering
  \begin{tikzpicture}[x=1cm, y=1cm, font=\scriptsize,
      eq/.style={ntcaixa, text width=1.7cm, minimum height=9mm},
      med/.style={circle, draw=corCinza, fill=white, line width=.5pt,
                  inner sep=0pt, minimum size=3.1mm, font=\tiny}]
    % ------------------------------------------------ fronteira -------
    \begin{scope}[on background layer]
      \fill[corFundo, rounded corners=3pt] (0,-1.35) rectangle (11.85,1.55);
      \draw[corTinta, line width=.7pt, dash pattern=on 3pt off 2pt,
            rounded corners=3pt] (0,-1.35) rectangle (11.85,1.55);
    \end{scope}
    \node[ntrot, anchor=north west, text=corTinta] at (0.15,1.45)
      {\textbf{Fronteira do indicador}};

    % ---------------------------------------------- cadeia processual --
    \node[eq] (pa1) at (1.75,0.30) {Pré-aquecimento};
    \node[eq] (des) at (3.90,0.30) {Dessalgação};
    \node[eq] (pa2) at (6.05,0.30) {Pré-aquecimento};
    \node[eq, draw=corFG, line width=.7pt, fill=corFG!7] (for) at (8.20,0.30)
      {Forno atmosférico};
    \node[ntcaixa, text width=1.6cm, minimum height=17mm,
          draw=corFG, line width=.7pt, fill=corFG!7] (col) at (10.70,0.20)
      {Coluna de destilação};

    \node[ntmini, anchor=north west] at (0.12,0.22) {crude};
    \draw[ntsetaf] (0.15,0.30) -- (pa1.west);
    \draw[ntsetaf] (pa1) -- (des);
    \draw[ntsetaf] (des) -- (pa2);
    \draw[ntsetaf] (pa2) -- (for);
    \draw[ntsetaf] (for) -- (col.west|-for);

    % ----------------------------------------------------- produtos ----
    \foreach \y/\t in {0.90/{gases e nafta}, 0.20/{petróleo e gasóleos},
                       -0.50/{resíduo atmosférico}} {
      \draw[ntsetaf] (11.50,\y) -- (12.15,\y);
      \node[ntmini, anchor=west, text width=2.5cm, align=left] at (12.20,\y) {\t};
    }

    % ------------------------------------------- entradas de energia ---
    % fuel gas: entra no forno, que e onde e queimado
    \draw[-{Stealth[length=4pt,width=3pt]}, draw=corFG, line width=.9pt]
      (8.20,-2.02) -- (for.south);
    \node[ntmini, anchor=north, text=corFG] at (8.20,-2.14) {fuel gás};
    \node[med] at (8.20,-1.35) {M};
    % eletricidade e vapores de alta: atravessam a fronteira, distribuidos
    % por varios equipamentos; a seta para na fronteira, sem atribuir destino
    \foreach \x/\c/\t in {1.75/corEE/{eletricidade}, 3.90/corV24/{vapor 24 e 10 bar}} {
      \draw[-{Stealth[length=4pt,width=3pt]}, draw=\c, line width=.9pt]
        (\x,-2.02) -- (\x,-1.20);
      \node[ntmini, anchor=north, text=\c] at (\x,-2.14) {\t};
      \node[med] at (\x,-1.35) {M};
    }
    \node[ntmini, anchor=north west, align=left] at (0.05,-2.50)
      {distribuídos por vários equipamentos; a eletricidade é medida ao mês e
       alocada ao dia};
    \draw[-{Stealth[length=4pt,width=3pt]}, draw=corV3, line width=.9pt]
      (10.70,-2.02) -- (col.south);
    \node[ntmini, anchor=north, text=corV3] at (10.70,-2.14) {vapor 3 bar};
    \node[med] at (10.70,-1.35) {M};

    % -------------------------------------- retornos e producao -------
    \draw[{Stealth[length=4pt,width=3pt]}-, draw=corV10, line width=.9pt]
      (6.05,2.00) -- (6.05,1.20);
    \node[ntmini, anchor=south, text=corV10] at (6.05,2.03)
      {vapor produzido na recuperação de calor};
    \node[med] at (6.05,1.55) {M};
    \draw[{Stealth[length=4pt,width=3pt]}-, draw=corCinza, line width=.7pt]
      (10.70,2.00) -- (col.north);
    \node[ntmini, anchor=south] at (10.70,2.03) {condensados};

    % ----------------------------------------------------- legenda -----
    \node[med] at (0.10,-3.20) {M};
    \node[ntmini, anchor=west] at (0.32,-3.20) {ponto de medição declarado};
    \draw[-{Stealth[length=4pt,width=3pt]}, draw=corCinza, line width=.9pt]
      (3.7,-3.20) -- ++(0.5,0);
    \node[ntmini, anchor=west] at (4.3,-3.20) {consumo bruto que entra};
    \draw[{Stealth[length=4pt,width=3pt]}-, draw=corV10, line width=.9pt]
      (7.6,-3.20) -- ++(0.5,0);
    \node[ntmini, anchor=west] at (8.2,-3.20) {produção creditada (abate ao bruto)};
  \end{tikzpicture}
  \caption[Fronteira energética da unidade de destilação atmosférica]%
  {Esquema funcional da unidade de destilação atmosférica e da fronteira
   que o indicador de desempenho usa. O desenho não tem escala nem caudais:
   representa a sequência de operações e os pontos onde entra e sai energia.
   Três noções de consumo convivem e não devem ser confundidas: o
   \emph{consumo bruto} é tudo o que atravessa a fronteira para dentro; o
   \emph{consumo líquido} desconta a produção creditada, como o vapor gerado
   na recuperação de calor; o \emph{consumo alocado} reparte por unidade
   aquilo que é medido num ponto comum, segundo uma regra declarada --- é
   o caso da eletricidade, medida ao mês para um conjunto de equipamentos e
   repartida por dia. Por isso, um desvio medido nesta fronteira pode ter
   origem no processo ou na regra de repartição.}
  \label{fig:fronteira-cdu}
\end{figure}


A refinaria de Sines é a única refinaria em operação em Portugal desde que a
Galp encerrou a refinação em Matosinhos, em 2021, e tem capacidade para
destilar 226 mil barris de crude por dia~\cite{Galp2026Industrial}. É certificada
segundo a ISO~50001 e tem quatro indicadores de desempenho energético
definidos: consumos específicos por unidade, índice de intensidade
energética, emissões de CO$_2$ associadas ao consumo e as retas que servem
de referência ao consumo de cada unidade (Secção~\ref{sec:contexto}).

\paragraph{A unidade.} A unidade de destilação atmosférica (CDU; CC na
designação interna) é a primeira por onde passa todo o crude. Segue o
esquema clássico: armazenagem, bateria de pré-aquecimento, dessalgação,
nova recuperação de calor, forno e coluna de destilação à pressão
atmosférica, com vapor de \emph{stripping} no fundo e nas colunas laterais,
refluxos circulantes que devolvem calor à bateria e cortes laterais que
seguem para as unidades a jusante~\cite{Gary2007}. Foi escolhida por três
razões que o Capítulo~\ref{cha:introducao} desenvolve: é, pela sua dimensão,
um uso significativo de energia por si só; é nela que o mecanismo de
incrustação da bateria de pré-aquecimento torna legível a falha de uma
referência estática; e quatro dos seus cinco vetores energéticos existem
com resolução diária.

\paragraph{A fronteira energética e os cinco vetores.} O consumo da CDU é
contabilizado em cinco vetores. O fuel gás alimenta o forno e é medido no
historiador de processo, com o poder calorífico do dia. O vapor entra em três
níveis de pressão, de 24, 10 e 3~bar; a unidade também produz vapor, e o vetor
de 10~bar é contabilizado como consumo líquido, que pode ser negativo. A eletricidade,
consumida sobretudo em bombas e aeroarrefecedores, só existe como total
mensal por unidade, distribuído uniformemente pelos dias. Os vetores são
convertidos numa unidade comum, a tonelada de gás natural equivalente por
dia (t~GNE/d), que é a convenção de contabilização em que a refinaria
orçamenta e reporta: soma energias de natureza diferente por equivalência
com um combustível de referência, e não é uma medida de energia primária nem
de exergia. As regras de conversão estão na Secção~\ref{sec:metricas}.

\paragraph{O índice de intensidade energética.} Ao lado das retas internas, a
refinaria acompanha o índice de intensidade energética da metodologia de
\emph{benchmarking} da indústria, que compara o consumo real com uma energia
padrão função da carga e das características de cada unidade. Serve para
comparar refinarias entre si e com os quartis do setor. Não é uma
referência de acompanhamento diário, e fica fora do diagnóstico
(Secção~\ref{sec:metricas}).

\paragraph{A referência em produção.} Para cada unidade, vetor e ano civil,
o modelo de dados estima por mínimos quadrados a reta $E = a\,Q + b$ entre o
consumo diário em t~GNE/d e a carga diária, usando só os dias em que a carga
atinge o limiar declarado para a unidade --- na CDU, $18\,600$~t/d, que deixa de
fora 10,4\% dos dias. O desvio-padrão dos resíduos do ano de referência
define bandas a $\pm1$, $\pm2$ e $\pm3\sigma$; a saída da banda de $\pm2\sigma$
é o alarme; um gráfico de somas acumuladas e oito regras de carta de controlo
completam a leitura (Secção~\ref{sec:metricas}). Cada ano tem a sua reta, e o
ano de referência pode ser escolhido. É, na tipologia da ISO~50006, um modelo
estatístico com uma variável explicativa (Tabela~\ref{tab:tipologia-enpi}). A
forma do modelo e a revisão periódica não são excluídas por nenhuma das
cláusulas analisadas na Secção~\ref{sec:enb}; a questão que a dissertação
coloca é outra: se os pressupostos de que dependem se verificam nesta unidade.

Os pressupostos que a Secção~\ref{sec:assimetria-matriz} extrai das normas e
os mecanismos da Secção~\ref{sec:processos-continuos} encontram-se aqui. A
carga é a única variável explicativa, e a incrustação, o crude e o estado da
rede de vapor ficam de fora; um ano civil é o período de treino; as bandas
assumem resíduos com escala conhecida; e a reestimação anual é a regra de
revisão. É este o objeto que o resto da dissertação examina: replicado de
forma exata a partir das tabelas de configuração oficiais
(Secção~\ref{sec:replicacao}), antes de qualquer crítica.

% [REV-2026-09-17][C2-08][P2][CONFIDENCIALIDADE]
\revisaotese{C2-08}{P2}{CONFIDENCIALIDADE}
{Confirmar o nível de detalhe da CDU com a orientação industrial}
{As Secções 2.4 e 2.5 descrevem a CDU ao nível de um esquema genérico de
destilação atmosférica e usam só números já presentes nos Caps. 4--6
(limiar de 18\,600~t/d, cinco vetores) e a capacidade pública da refinaria.
Confirmar que a orientação industrial aceita este nível de detalhe e a
descrição da configuração (dessalgação, refluxos circulantes, produção de
vapor de 10~bar); não foi usado nenhum documento de licenciador.}%
% [/REV-2026-09-17]

%% ====================================================================
%% PENDENCIAS DESTE CAPITULO (2026-09-19; retiradas do corpo, onde eram
%% \todo{} e saiam impressas na margem do PDF, cortadas pelo rebordo):
%%
%% 1. Confirmar com a orientacao industrial qual a edicao da ISO 50006
%%    que enquadra o SGE da refinaria (a analise usa a de 2014; existe
%%    uma segunda edicao de 2023, nao consultada). Sec. 2.3.3, "Nota
%%    sobre as edicoes analisadas".
%% 2. Confirmar a formulacao do limiar R2 >= 0,75 na edicao corrente do
%%    IPMVP (EVO 10000-1). Por agora o texto e a Tabela 2.3 dao-no como
%%    relatado por Ruiz e Bandera, sem consulta do protocolo -- o que ja
%%    esta declarado e e defensavel como esta. Sec. 2.3.3.
%% ====================================================================
